\begin{figure}[htbp]
\centering
\begin{tikzpicture}[x=1cm,y=1cm,>=Latex,line cap=round,line join=round,font=\sansmath\sffamily]
  \tikzset{
    axis/.style={->,line width=1pt,black},
    basegrid/.style={gray!55,line width=1pt,dash pattern=on 4pt off 3pt},
    movinggrid/.style={magenta!28,line width=1pt,dash pattern=on 6pt off 3pt},
    worldcar/.style={violet!90!black,line width=1pt},
    worldball/.style={blue!90!black,line width=1pt}
  }

  \def\xmin{-2}
  \def\xmax{6}
  \def\ymax{5}

  % background square grid
  \draw[basegrid] (\xmin,0) grid (\xmax,\ymax);

  % moving-frame oblique grid: all lines are parallel to the car worldline
  % and clipped to stay strictly inside the square grid frame.
  \begin{scope}
    \clip (\xmin,0) rectangle (\xmax,\ymax);
    % 2 lines on the left, 5 lines on the right of the main purple line
    \foreach \x in {-1,0,1,2,3,4,5} {
      \draw[movinggrid] (\x,0) -- (\x+5,5);
    }
    \foreach \y in {0,1,2,3,4} {
      \draw[movinggrid] (\xmin,\y) -- (\xmin+5,\y+5);
    }
  \end{scope}

  % axes
  \draw[axis] (\xmin,0) -- (7,0) node[below right=2pt] {$x\equiv x'$};
  \draw[axis] (0,0) -- (0,6) node[above left=1pt] {$t$};

  % worldlines
  \draw[worldcar] (0,0) -- (5,5);
  \draw[axis] (5,5) -- (6,6) node[above right=1pt] {$t'$};
  \draw[worldball] (1,0) -- (3,5);

  % origin mark
  \fill[red!85!black] (0,0) circle (2pt);
  \node[below left=1pt,text=red!85!black] at (0,0) {$O$};

  % distance annotation: one square = 10 m
  \draw[gray!75,<->,line width=1pt] (0,-1) -- (1,-1);
  \draw[gray!75,line width=1pt] (0,-1) -- (0,0);
  \draw[gray!75,line width=1pt] (1,-1) -- (1,0);
  \node[below=2pt,text=red!85!black] at (0.5,-1) {$10\,\mathrm{m}$};
\end{tikzpicture}
\caption{Minh họa sự chồng chập xét trong hệ quy chiếu cây}
\label{P1_tikz2}
\end{figure}